\documentclass[a4paper,twoside]{article}

\usepackage{morewrites}

\usepackage[svgnames,dvipsnames,x11names]{xcolor}
\usepackage{graphicx}
\usepackage[english]{babel}
\usepackage[colorlinks=true, linkcolor=NavyBlue, urlcolor=NavyBlue, citecolor=NavyBlue]{hyperref}
\usepackage[a4paper]{geometry}
\usepackage{ts-typesetting}
\usepackage{ts-fonts}
\usepackage{xcolor}
\usepackage{epigraph}
\usepackage{marginnote}
\usepackage{multicol}
\usepackage{progressbar}
\usepackage{graphicx}
\usepackage{tcolorbox}
\usepackage{fvextra}
\usepackage{amsmath}
\usepackage{float}
\usepackage{seqsplit}
\renewcommand*{\marginfont}{
\footnotesize
\sloppy
\emergencystretch=1em
\hyphenpenalty=50
\exhyphenpenalty=50
}
\newlength{\tsmarginparavail}
\newlength{\tsmarginparalt}
\newlength{\tsmarginparbuf}
\setlength{\tsmarginparbuf}{6mm}
\makeatletter
\AtBeginDocument{
\setlength{\tsmarginparavail}{\dimexpr\paperwidth-\textwidth-1in-\oddsidemargin-\marginparsep-\tsmarginparbuf\relax}
\ifdim\tsmarginparavail<\z@\setlength{\tsmarginparavail}{\z@}\fi
\if@twoside
\setlength{\tsmarginparalt}{\dimexpr1in+\evensidemargin-\marginparsep-\tsmarginparbuf\relax}
\ifdim\tsmarginparalt<\z@\setlength{\tsmarginparalt}{\z@}\fi
\ifdim\tsmarginparalt<\tsmarginparavail
\setlength{\tsmarginparavail}{\tsmarginparalt}
\fi
\fi
\ifdim\tsmarginparavail<\marginparwidth
\setlength{\marginparwidth}{\tsmarginparavail}
\fi
}
\makeatother
\usepackage{ts-code}

\usepackage{lastpage}
\usepackage{refcount}
\makeatletter
\def\ps@plain{
\let\@oddhead\@empty
\let\@evenhead\@empty
\def\@oddfoot{
\hfil
\ifnum\getpagerefnumber{LastPage}>1\relax
\thepage
\fi
\hfil}
\let\@evenfoot\@oddfoot
}
\makeatother

\title{Bibliography}
\author{}\date{}
\begin{document}
\maketitle
\section{Add Bibliography Entries}
To include bibliography entries in your document, provide one or more \tslogo{BibTeX} files as inputs when running TeXSmith. TeXSmith will parse these files and integrate the references into your document during the conversion process. For example:
\begin{tscode}[lang=bash, engine=pygments]
\begin{Verbatim}[commandchars=\\\{\},breaklines, breakanywhere, commandchars=\\\{\}]
texsmith\PY{+w}{ }paper.md\PY{+w}{ }nature.bib\PY{+w}{ }ieee.bib\PY{+w}{ }\PYZhy{}o\PY{+w}{ }paper.pdf
\end{Verbatim}
\end{tscode}
\section{List Bibliography Entries}
Use the \tscodeinline{-\allowbreak{}-\allowbreak{}list-\allowbreak{}bibliography} flag to inspect \tslogo{BibTeX} files before running a conversion or build. It helps catch parsing issues, duplicate entries, and empty datasets early in your workflow. For example in the \tscodeinline{paper} example project:
\begin{tscode}[lang=text, stretch=0.5, engine=pygments]
\begin{Verbatim}[commandchars=\\\{\},breaklines, breakanywhere, commandchars=\\\{\}]
\PYZdl{} texsmith cheese.md cheese.bib \PYZhy{}\PYZhy{}list\PYZhy{}bibliography

                    Bibliography Files
┌──────────────────────────────────────────────┬─────────┐
│ File                                         │ Entries │
├──────────────────────────────────────────────┼─────────┤
│ /home/ycr/texsmith/examples/paper/cheese.bib │       2 │
│ Total                                        │       2 │
└──────────────────────────────────────────────┴─────────┘
                              Jaoac2019 (article)
  Title    Determination of Moisture in Cheese and Cheese Products
  Year     2019
  Journal  Journal of AOAC INTERNATIONAL
  Authors  Jr Bradley Robert L, Margaret A Vanderwarn
  Sources  /home/ycr/texsmith/examples/paper/cheese.bib
  Abstract Variables related to oven\PYZhy{}drying samples of cheese and cheese
           products to determine moisture...
  Doi      10.1093/jaoac/84.2.570
  Eprint   https://academic.oup.com/jaoac/article\PYZhy{}pdf/84/2/570/32415847/jaoac…
  Issn     1060\PYZhy{}3271
  Month    11
  Number   2
  Pages    570\PYZhy{}592
  Url      https://doi.org/10.1093/jaoac/84.2.570
  Volume   84

                             Prentice1993 (inbook)
  Title     Cheese Rheology
  Year      1993
  Authors   J. H. Prentice, K. R. Langley, R. J. Marshall
  Sources   /home/ycr/texsmith/examples/paper/cheese.bib
  Abstract  Rheology is formally defined as the study of the flow and
            deformation of matter...
  Address   Boston, MA
  Booktitle Cheese: Chemistry, Physics and Microbiology: Volume 1 General
            Aspects
  Doi       10.1007/978\PYZhy{}1\PYZhy{}4615\PYZhy{}2650\PYZhy{}6\PYZus{}8
  Isbn      978\PYZhy{}1\PYZhy{}4615\PYZhy{}2650\PYZhy{}6
  Pages     303\PYZhy{}\PYZhy{}340
  Publisher Springer US
  Url       https://doi.org/10.1007/978\PYZhy{}1\PYZhy{}4615\PYZhy{}2650\PYZhy{}6\PYZus{}8

    Bibliography Summary
┌───────────────────┬───────┐
│ Category          │ Count │
├───────────────────┼───────┤
│ Total entries     │     2 │
│ From cheese.bib   │     2 │
│ From front matter │     0 │
│ From DOI fetches  │     0 │
└───────────────────┴───────┘
\end{Verbatim}
\end{tscode}
\section{Behaviour}
TeXSmith loads every provided \tscodeinline{.bib} file using \tscodeinline{pybtex} and analyzes their contents, checks for duplicate keys, and verifies parsing integrity. It emits warnings for any issues found and prints a summary table of the number of entries per file.

The \tscodeinline{-\allowbreak{}-\allowbreak{}list-\allowbreak{}bibliography} flag has the following behavior:
\begin{itemize}
\item Prints a formatted table summarising the number of entries per file.
\item Emits warnings for files that fail to parse, contain duplicate keys, or are empty.
\item Highlights issues detected by TeXSmith's bibliography loader (e.g. conflicting entries sourced from multiple files).
\item Exits before rendering anything else, so you can run it as a fast preflight step.
\end{itemize}
\end{document}
